本文面向树莓派 5 这类无 GPU ARM 边缘平台上的 LeRobot 在线控制任务，
围绕 SO-101 机械臂、双摄像头输入和 ACT 策略模型，完成了对
“观测--推理--执行”闭环的系统级性能剖析。通过将核心控制循环中的阶段计时、
系统调用追踪、内核阻塞分析和硬件通信链路建模结合起来，
本文建立了一套覆盖 AI 推理框架、Python 运行时、Linux 内核与硬件 I/O 的跨层分析方法。
该方法能够把端到端 FPS 下降进一步定位到具体软件层次和等待来源，
避免只根据单点计时结果判断系统瓶颈。

实验结果表明，当前平台的主要性能瓶颈并不在摄像头读取、舵机写入或单个 Linux 系统调用上，
而在 ARM CPU 上执行 ACT 模型前向推理的绝对耗时。摄像头路径已经通过 OpenCV 后台线程实现异步预取，
主线程观测开销较低；SCServo 舵机读写的毫秒级延迟主要受 1 Mbps 半双工总线和协议交互约束，
系统调用本身只占很小比例。相比之下，单次完整 ACT 推理达到秒级耗时，
其中 ResNet18 视觉编码和 Transformer 模块占据主要比例，
是限制在线控制频率和造成同步执行 FPS 不足的核心原因。

基于上述结论，本文进一步验证了推理异步化对机器人在线控制的实际收益。
在 chunk\_size 为 100、目标 30 FPS 的实验设置下，
异步推理触发阈值取 0.30 时有效控制频率达到 27.42 FPS，
相对同步 baseline 提升 70.0\%，并在稳态阶段显著减少动作队列 stall。
这说明在不改变模型结构和硬件平台的条件下，
通过基础软件层面的流水线调度可以有效回收动作执行期间的空窗，
缓解长耗时推理对控制闭环的阻塞。

综上，本文证明了边缘机器人系统中的性能问题不能简单归因于某个模型、
某个外设或某次系统调用，而需要放在完整的软件栈和控制闭环中分析。
面向 LeRobot 这类机器人学习框架，系统级分层 profiling 不仅能够解释当前瓶颈的来源，
也能为后续推理运行时优化、操作系统调度改进和机器人基础软件架构设计提供依据。
